\documentclass[a4paper]{scrartcl}
\usepackage[utf8]{inputenc}
\usepackage[english]{babel}
\usepackage{graphicx}
% SVG support: converted to vector PDF by Inkscape at compile time (needs -shell-escape, see latexmkrc).
% inkscapelatex=false keeps the diagram's own text and fonts instead of re-typesetting them with LaTeX.
\usepackage[inkscapelatex=false]{svg}
\usepackage{lastpage}
\usepackage{pgf}
\usepackage{wrapfig}
\usepackage{fancyvrb}
\usepackage{listings}
\usepackage{xcolor}
\usepackage{url} % for the repository link in the Result chapter
% The two ER diagrams are too wide for a portrait text block. Each one gets a page of its
% own, with the diagram turned sideways and the margins reduced so it is as large as
% possible. lscape rotates only the content, so the page itself stays portrait (use
% pdflscape instead to also set /Rotate and have viewers show the page turned).
% geometry is loaded with pass so it changes nothing except to provide \newgeometry.
\usepackage{lscape}
\usepackage[pass]{geometry}
% \diagrampage{<file>}{<caption>}{<label>}: one rotated, nearly full-page figure.
\newcommand{\diagrampage}[3]{%
  \newgeometry{margin=10mm}%
  \begin{landscape}%
    \vspace*{\fill}
    % Inside landscape, \linewidth is the rotated (long) page direction, while the
    % available height is bounded by \paperwidth. keepaspectratio picks whichever
    % of the two limits binds, so one macro fits both a wide and a squarer diagram.
    \begin{figure}[h!]
      \centering
      \includesvg[width=\linewidth,height=0.78\paperwidth,keepaspectratio]{#1}
      \caption{#2}
      \label{#3}
    \end{figure}%
    \vspace*{\fill}
  \end{landscape}%
  \restoregeometry%
}
% Style for SQL / PL/pgSQL code listings
\lstdefinestyle{sql}{
  language=SQL,
  morekeywords={RETURNS,TRIGGER,DECLARE,BEGIN,END,IF,THEN,RAISE,EXCEPTION,RETURN,LANGUAGE,FUNCTION,REPLACE,EXECUTE,BEFORE,EACH,ROW,INTEGER,LIMIT,INTO,INNER,JOIN},
  basicstyle=\ttfamily\footnotesize,
  keywordstyle=\color{blue!70!black}\bfseries,
  commentstyle=\color{green!45!black}\itshape,
  stringstyle=\color{red!65!black},
  frame=single,
  breaklines=true,
  columns=fullflexible,
  keepspaces=true,
  showstringspaces=false
}
% Long identifiers such as UNIQUE(course_instance_id, teaching_activity_id) are single
% unbreakable words, so TeX pushed them past the right margin (one line stuck out by
% 100pt). Allow a line break after each underscore, without inserting a hyphen that
% could be mistaken for part of the name. emergencystretch mops up the tighter lines.
\renewcommand{\_}{\textunderscore\allowbreak}
\emergencystretch=2em

\usepackage{fancyhdr}
\pagestyle{fancy}

% Create header and footer
\headheight 27pt
\pagestyle{fancyplain}
\lhead{\footnotesize{Data Storage Paradigms, IV1351}}
\chead{\footnotesize{Project task 2 report}}
\rhead{}
\lfoot{}
\cfoot{\thepage\ (\pageref{LastPage})}
\rfoot{}

\title{Data Storage Paradigms, IV1351}
\subtitle{Project task 2 report}
\author{}
\date{\today}

\begin{document}

\maketitle
\noindent\textbf{Project members:} \\ \hfill
[Mo Wang, mowang@ug.kth.se] \\ \hfill
% [Enter name, email] \\ \hfill
% [Enter name, email] \\ \hfill

\section*{Declaration:}

By submitting this assignment, it is hereby declared that all group members listed above have contributed to the solution. It is also declared that all project members fully understand all parts of the final solution and can explain it upon request.

It is furthermore declared that the solution below is a contribution by the project members only, and specifically that no part of the solution has been copied from any other source (except for lecture slides at the course IV1351), no part of the solution has been provided by someone not listed as a project member above, and no part of the solution has been generated by a system.

% \section*{Tips for Report Writing}
% \textbf{REMOVE THIS SECTION BEFORE SUBMITTING THE REPORT.}\\

% \noindent \textit{The target audience has exactly the same skills as the author, except they do not know anything at all about the specific application described in the report.} \\

% Consider the following:

% \begin{itemize}
%   \item \textbf{The report must be \textit{centered around the requirements}. Which are they (Introduction), how did you work to meet them (Method), what is the solution that meets them (Result), and how can you be sure they are met (Discussion). This is the IMRaD method.} The requirements on the Introduction, Method, Result and Discussion chapters are described below under each chapter.

%   \item Is spelling and grammar correct? Is spoken language avoided?

%   \item Does the report have a good structure with sections, subsections and paragraphs?

%   \item Is the text clarified with images and/or other figures, and with links to the code in your Git repository? Remember that all figures (images, tables, graphs, code listings, etc) shall be numbered and have a short explaining text.
% \end{itemize}

\section{Introduction}
\label{sec:introduction}

% \textbf{This chapter tells \textit{what} are you going to do.} 
% Explain the task and the requirements on the solution. It's important to clearly state the requirements. \textit{Also specify which other student you worked with when solving the tasks, or if you worked alone.}

This report presents a set of analytical SQL queries for the course database project, which was already designed in project task 1. The project task was carried out alone. The database stores courses, course layout versions, course instances, planned teaching activities, employees and their teaching allocations at a university. The analytical SQL queries calculate planned hours, actually allocated hours per course instance and per teacher, the load of each teacher, the teachers who are allocated to many course instances in the same period and mismatch rate between actual allocated hours and planned hours.

Additionally, scripts that scale up the database and performance analysis using \texttt{EXPLAIN ANALYZE} are implemented as well. This lays the foundation of performance analysis of the database, as well as improving the performance using suitable indexes and materialized views for the most frequent queries, according to the list of query workload with approximate query frequencies.

\subsection{Requirements of project task 2}
\label{sec:requirements}

The requirements of project task 2 concern OLAP (Online Analytical Processing) queries, which are executed manually to analyze and summarize both the planned and allocated teaching load. The queries are run against the course database, which was created earlier in project task 1. 

\subsubsection{Queries}
There are four mandatory analytical queries and one higher grade query as requirement, as well as scaling up the database and optimizing the five given queries in the report. 

The five analytical queries calculate planned hours, allocated hours for each course instance, allocated hours for each teacher, number of course instances each teacher has and mismatch between allocated and planned hours.

For planned hours, the total planned hours of every course instance of the current year is calculated, including multiplication factors. Each row contains one course instance, with the planned hours broken down per teaching activity. 

The query that calculates allocated hours for a course instance, all allocation hours of one teaching activity in a specific course instance is summed up, including the multiplication factor. The result is listed per teacher and per teaching activity for a specific course instance.

Another similar query that calculates the allocated hours but for a teacher proceeds similar to the query that calculates the allocated hours for a course instance, but for a specific teacher instead of a course instance.

Finally, a query that lists the employment IDs and names of all teachers allocated to more than a specific number of course instances in the current period, together with that number, is also required.

\subsubsection{Scaling up database and performance improvements}

For higher grades, the efficiency of at least one query must be analysed with \texttt{EXPLAIN ANALYZE}, and the database must contain enough data to show that all queries work as intended. For the higher grade part, in order to observe noticeable performance improvement differences, the database must be scaled up, so that the largest table contains at least 1,000,000 rows.

Besides, there is another query that should list course instances that have mismatches between allocated and planned hours with more than 15\%. A wrong version of this query, that gives an incorrect result, must also be written and explained.

According to the workload below with approximate number of query execution each day,
\begin{itemize}
  \item planned hours per course instance (8 times per day)
  \item actual allocated hours per teacher per course instance (12 times per day)
  \item teacher load per period (5 times per day)
  \item course instances where planned and allocated hours differ by more than 15\% (3 times per day)
  \item teachers allocated to more than N course instances in the current period (20 times per day)
\end{itemize}

it is required that suitable indexes and materialized views are added where they are needed with explanation for each choice. In order to observe any noticeable difference, \texttt{EXPLAIN ANALYZE} is run before and after. Additionally, one index that PostgreSQL does not use is found and explained.

Also, it is shown that a materialized view is out of date after a new allocation is inserted, together with the time it takes to refresh it.

Whether indexes and materialized views are reasonably justified in terms of performance gain, and what would change if the workload changed, are explained and analyzed in the discussion as well. 


\section{Comments About the Course}

% \textbf{This section is optional, but please at least write approximately how much time you spent on the assignment}, including lectures, labs, tutorials and seminars. This is of great help for course evaluation.

% Also, any other comment(s) related to this course offering or to coming offerings is much appreciated. 

\section{References}
\label{sec:references}

% thebibliography is used instead of BibTeX, so that the report compiles with pdflatex
% alone and the reference list is kept in this single source file. Cite with \cite{key}.
\renewcommand{\refname}{}
\vspace{-2\baselineskip}
\begin{thebibliography}{9}

\bibitem{elmasri}
Ramez Elmasri and Shamkant B. Navathe.
\textit{Fundamentals of Database Systems}, 7th edition.
Pearson, 2016.

\bibitem{silberschatz}
Abraham Silberschatz, Henry F. Korth and S. Sudarshan.
\textit{Database System Concepts}, 7th edition (ISE).
McGraw-Hill, 2019.

\bibitem{project}
Leif Lindbäck and the IV1351 teaching staff.
\textit{Project: Course Layout and Teaching Load Allocations}.
Course page, Data Storage Paradigms IV1351, KTH, HT26.

\bibitem{task1}
\textit{Project Task 1, Logical and Physical Model}.
Course page, Data Storage Paradigms IV1351, KTH, HT26.

\bibitem{task2}
\textit{Project Task 2, SQL}.
Course page, Data Storage Paradigms IV1351, KTH, HT26.

\bibitem{tips}
\textit{Tips and Tricks for Project Task 1} (\texttt{seminar1-tips-and-tricks.pdf}).
Course material, Data Storage Paradigms IV1351, KTH, HT26.

\bibitem{tips2}
\textit{Tips and Tricks for Project Task 2} (\texttt{seminar2-tips-and-tricks.pdf}).
Course material, Data Storage Paradigms IV1351, KTH, HT26.

\bibitem{criteria}
\textit{Assessment Criteria For Seminar 1, Logical and Physical Model}
(\texttt{assessment-criteria-seminar1.pdf}).
Course material, Data Storage Paradigms IV1351, KTH, HT26.

\bibitem{criteria2}
\textit{Assessment Criteria For Seminar 2} (\texttt{seminar2-assessment-criteria.pdf}).
Course material, Data Storage Paradigms IV1351, KTH, HT26.

\bibitem{conceptuallecture}
Leif Lindbäck.
\textit{The Conceptual Model}. Recorded lecture series, about three hours.
Data Storage Paradigms IV1351, KTH, HT26.

\bibitem{logicallecture}
Leif Lindbäck.
\textit{Logical and Physical Models}. Recorded lecture series.
Data Storage Paradigms IV1351, KTH, HT26.

\bibitem{normalisationlecture}
\textit{Normalisation}. Lecture 4, Data Storage Paradigms IV1351, KTH, HT26.

\bibitem{sqllecture}
\textit{SQL, the Structured Query Language}. Lecture 3, Data Storage Paradigms IV1351, KTH, HT26.

\bibitem{postgresdocs}
The PostgreSQL Global Development Group.
\textit{PostgreSQL 16 Documentation}.
\url{https://www.postgresql.org/docs/16/}.

\end{thebibliography}


\end{document}
